% Λύση του 4ου προβλήματος της ενότητας «Άλυτα Προβλήματα».
\subsection*{Λύση Προβλήματος 4 --- Ελαστικότητα φορολογητέου εισοδήματος}

\begin{alist}
  \item Για $\alpha=1{,}5=\dfrac32$ και
        $\varepsilon=0{,}25=\dfrac14$ έχουμε
        \[
          \alpha\varepsilon
          =\frac32\cdot\frac14
          =\frac38.
        \]
        Επομένως, ο κανονικοποιημένος δείκτης εσόδων είναι
        \[
          {R(x)=x(1-x)^{3/8}},
          \qquad x\in[0,1].
        \]

  \item Για $x\in(0,1)$, εφαρμόζουμε τον κανόνα παραγώγισης γινομένου
        και τον κανόνα της αλυσίδας:
        \begin{align*}
          R'(x)
          &=x'(1-x)^{\alpha\varepsilon}
            +x\left[(1-x)^{\alpha\varepsilon}\right]'\\
          &=(1-x)^{\alpha\varepsilon}
            -\alpha\varepsilon x
             (1-x)^{\alpha\varepsilon-1}\\
          &=(1-x)^{\alpha\varepsilon-1}
            \left[(1-x)-\alpha\varepsilon x\right]\\
          &=(1-x)^{\alpha\varepsilon-1}
            \left[1-(1+\alpha\varepsilon)x\right].
        \end{align*}
        Άρα,
        \[
          {
          R'(x)=(1-x)^{\alpha\varepsilon-1}
          \left[1-(1+\alpha\varepsilon)x\right]}.
        \]

  \item Για $x\in(0,1)$ ισχύει
        \[
          (1-x)^{\alpha\varepsilon-1}>0.
        \]
        Επομένως, το πρόσημο της $R'(x)$ καθορίζεται αποκλειστικά από
        τον παράγοντα
        \[
          1-(1+\alpha\varepsilon)x.
        \]
        Έχουμε
        \[
          R'(x)=0
          \Longleftrightarrow
          x=\frac{1}{1+\alpha\varepsilon}=x^*.
        \]
        Επίσης,
        \[
        \begin{array}{c|ccc}
          x & (0,x^*) & x^* & (x^*,1)\\ \hline
          R'(x) & + & 0 & -\\
          R(x) & \nearrow & \text{μέγιστο} & \searrow
        \end{array}
        \]
        Άρα η $R$ είναι γνησίως αύξουσα στο $[0,x^*]$ και γνησίως
        φθίνουσα στο $[x^*,1]$. Επειδή
        \[
          R(0)=R(1)=0,
        \]
        η συνάρτηση παρουσιάζει ολικό μέγιστο στο
        \[
          {x^*=\frac{1}{1+\alpha\varepsilon}}.
        \]

  \item Για $\alpha\varepsilon=\dfrac38$ βρίσκουμε
        \[
          x^*
          =\frac{1}{1+\frac38}
          =\frac{1}{\frac{11}{8}}
          =\frac{8}{11}
          \simeq0{,}727.
        \]
        Συνεπώς, στο συγκεκριμένο μοντέλο,
        \[
          {x^*\simeq72{,}7\%}.
        \]

        Το αποτέλεσμα αφορά τον φοροεισπρακτικά βέλτιστο ανώτατο
        οριακό συντελεστή κάτω από τις συγκεκριμένες υποθέσεις και τις
        συγκεκριμένες τιμές των παραμέτρων. Δεν αποτελεί καθολική
        πρόταση για όλους τους φόρους ή όλες τις οικονομίες. Η τιμή
        μεταβάλλεται όταν αλλάζουν η κατανομή των εισοδημάτων, η
        ελαστικότητα της φορολογητέας βάσης, η φοροδιαφυγή, το
        φορολογικό σύστημα ή οι κοινωνικοί στόχοι της πολιτικής.
\end{alist}
